\section{Statistical model and the meaning of a certificate}
\label{sec:foundations}
A measurement schedule determines which responses enter the likelihood. With
correlated errors, changing the schedule changes the covariance matrix that
must be inverted. This section specifies that statistical target, compares it
quantitatively with a commonly used surrogate, and defines what a global
certificate establishes. The distinctions apply both to sensor installation
and to sampling-time selection.

\subsection{Selected observations and local information}
\label{subsec:model}
There are $n$ candidate measurement packets. Packet $t$ contains $d_t$ scalar
responses, acquired together when that packet is selected. Set
$d_{\mathrm{tot}}=\sum_t d_t$; scalar candidates have $d_t=1$. The candidate
observation vector satisfies
\begin{equation}
 Y=\mu(\theta)+\epsilon,\qquad \epsilon\sim N(0,R),\qquad
 R\in\Splusplus^{d_{\mathrm{tot}}},
 \label{eq:observation-model}
\end{equation}
where $\theta\in\R^p$, the mean is differentiable at a specified nominal
parameter $\theta_0$, and the known covariance is independent of $\theta$.
A schedule is chosen before the data are observed. It determines which
coordinates are observed and does not alter the process or covariance.
Adaptive sampling and interventions on the process are outside this model.

Write $F=\partial\mu/\partial\theta|_{\theta_0}\in
\R^{d_{\mathrm{tot}}\times p}$. For a schedule $S\subseteq\{1,\ldots,n\}$,
$E_S$ selects its scalar coordinates, $F_S=E_SF$, and
$R_{SS}=E_SRE_S^\top$. Thus a matrix subscript $S$ includes every coordinate
of its selected packets. The feasible family $\calF$ is finite and nonempty;
it may encode cardinality, costs, installation, and spacing constraints.
Complexity claims require the more specific families stated in each theorem;
an arbitrary feasible family is not assumed to have an efficient oracle.

Marginalizing unobserved responses gives
$Y_S\sim N(\mu_S(\theta),R_{SS})$. Its score at $\theta_0$ and its Fisher
information are
\begin{align}
 \nabla_\theta\log p_\theta(Y_S)|_{\theta_0}
 &=F_S^\top R_{SS}^{-1}(Y_S-\mu_S(\theta_0)),\\
 I(S)&=F_S^\top R_{SS}^{-1}F_S.
 \label{eq:true-information}
\end{align}
The second identity is the second moment of the first. It is the exact
Fisher information of the stated Gaussian model, also for a nonlinear mean
evaluated at $\theta_0$. The approximation in local design lies in using a
nominal parameter for the unknown one, or local information to predict
estimation accuracy; it does not justify interchanging selection and
inversion. This selected-covariance formulation is established in correlated
design and sensor selection \citep{Patan2007,Liu2016,Pazman2022}.

We optimize the augmented matrix
\begin{equation}
 J(S)=J_0+I(S),\qquad J_0\in\Splus^p,
 \label{eq:augmented-information}
\end{equation}
with $I(\varnothing)=0$. A fixed $J_0$ can represent independent prior
information or a declared regularization; these interpretations must be
reported separately. In a linear Gaussian model, a Gaussian prior with
precision $J_0\succ0$ makes $J(S)^{-1}$ the posterior covariance. For a
nonlinear mean, adding a matrix at one nominal parameter does not generally
compute exact posterior covariance or Bayesian expected information.
A numerical diagonal regularization changes the design problem and does not
establish information supplied by an actual prior experiment.

The fixed-covariance assumption is substantive. If $R=R(\theta)$, the
$(a,b)$ entry of Fisher information includes the additional term
\begin{equation}
 \frac12\tr\!\left(R_{SS}^{-1}
 \frac{\partial R_{SS}}{\partial\theta_a}
 R_{SS}^{-1}\frac{\partial R_{SS}}{\partial\theta_b}\right).
 \label{eq:covariance-information}
\end{equation}
Patan and Bogacka explicitly consider this dependence \citep{Patan2007}.
Our results assume a fixed covariance and do not apply to a model that
requires this term.
For matrices, $A\preceq B$ means $B-A$ is positive semidefinite,
$\ip{A}{B}=\tr(A^\top B)$, and $\norm{A}_2$ is the spectral norm.
A singular prior is allowed; a singular observation covariance is not allowed
in \eqref{eq:observation-model}. A singular Gaussian likelihood requires
analysis on its support, possibly with parameter-dependent deterministic
constraints, and cannot generally be handled by replacing inverses with
pseudoinverses.

\subsection{Criteria, rank, and parameter coordinates}
\label{subsec:criteria}
We use the following conventional criteria \citep{Sagnol2015,Hainy2025}:
\begin{center}
\begin{tabular}{lll}
\toprule
Criterion & Objective & Domain\\
\midrule
D-optimality & maximize $\log\det J(S)$ & $J(S)\succ0$\\
Weighted trace information & maximize $\tr(WJ(S))$, $W\succeq0$ & $J(S)\succeq0$\\
A-optimality & minimize $\tr(J(S)^{-1})$ & $J(S)\succ0$\\
E-optimality & maximize $\lambda_{\min}(J(S))$ & $J(S)\succeq0$\\
\bottomrule
\end{tabular}
\end{center}
Trace information is distinct from trace inverse information; we use the
descriptive name ``trace information'' to avoid ambiguity with other uses
of the letter T in design theory. The criterion called A-optimality in
\citet{Wang2024} maximizes $\tr J$, whereas our A-optimality minimizes
$\tr(J^{-1})$.

All logarithms are natural. We may extend $\log\det J=-\infty$ to singular
positive semidefinite matrices. Every finite D-optimality certificate must
include a positive definite feasible information matrix. If none exists,
the optimum determinant is zero and there is no finite full-parameter
log-determinant optimum. Inverse-based criteria are used only when their
inverses exist. In a linear model with no prior, a contrast $c^\top\theta$
is estimable exactly when $c\in\range I(S)$, and its minimum unbiased
variance is $c^\top I(S)^\dagger c$. An analogous augmented-matrix expression
is an algebraic or prior-regularized criterion and requires its declared
interpretation. A restriction of all information matrices is valid for a
declared reduced parameter model, or on a common supported invariant range
expressed in an orthonormal basis. If complementary parameters remain unknown
nuisance parameters, use the contrast covariance accounting for those
parameters (in a nonsingular block model, by inverting the nuisance-adjusted
Schur complement). Mere compression followed by inversion need not give that
covariance.
A pseudodeterminant on design-dependent ranges would define a different
objective.

For positive definite information, maximizing $\log\det J$, $\det J$, and
$(\det J)^{1/p}$ gives the same schedules. The log-determinant and determinant
root are concave in $J$; the determinant itself is generally not concave for
$p>1$. Under an invertible change $\theta=T\eta$, sensitivities become $FT$
and the prior must become $T^\top J_0T$. The log-determinant changes by the
schedule-independent constant $2\log|\det T|$. D-optimal schedules and
relative D-efficiencies are therefore invariant. Unweighted trace information,
A-optimality, and E-optimality are generally coordinate-dependent. A weight
or contrast can be transformed with the coordinates when its interpretation
is to be preserved.

Positive definite local information does not imply global identifiability.
Observing only the intermediate species in consecutive first-order reactions
gives the mean
\begin{equation}
 B(t;A_0,k_1,k_2)=\frac{A_0k_1}{k_2-k_1}
 (e^{-k_1t}-e^{-k_2t}),\qquad A_0,k_1,k_2>0,\quad k_1\ne k_2.
 \label{eq:rate-swap}
\end{equation}
Substitution shows that $(A_0,k_1,k_2)$ and
$(A_0k_1/k_2,k_2,k_1)$ give the same entire trajectory. With fixed covariance,
they give the same likelihood for every schedule. Nevertheless, at
$(A_0,k_1,k_2)=(1,1,2)$, the sensitivity matrix with rows at times
$\log2$, $2\log2$, $3\log2$ and columns in the order $(A_0,k_1,k_2)$
has determinant $-(\log2)^2/2048$. It has full rank, so any positive definite
observation covariance gives positive definite local data information for
these three observations without a prior. A design certificate at one nominal
parameter, including a regularized certificate, does not resolve
this ambiguity. A physical restriction such as a known rate ordering changes
the admissible parameter model and must be stated when used.

\subsection{Selection before inversion and full-inverse gating}
\label{subsec:gating}
An alternative expression weights precomputed full-precision entries by
pairs of selected coordinates. With $D_S=E_S^\top E_S$, it is
\begin{equation}
 J_{\gate}(S)=J_0+F^\top D_SR^{-1}D_SF
 =J_0+F_S^\top(R^{-1})_{SS}F_S.
 \label{eq:gated-information}
\end{equation}
Liu et al. distinguish this expression from marginal information and analyze
its weak-correlation approximation \citep[Sec.~IV]{Liu2016}. The distinction
is theirs; we add proofs and quantitative consequences for interpreting
optimization results.

\begin{proposition}[Information inflation and equality]
\label{prop:gating}
For every schedule, $J(S)\preceq J_{\gate}(S)$. For a nonempty proper
selection, partition the scalar coordinates into selected and unselected
sets and write
\[
 R=\begin{pmatrix}A&B\\B^\top&C\end{pmatrix},\qquad
 A=R_{SS},\quad C=R_{S^cS^c}.
\]
Then
\begin{equation}
 J_{\gate}(S)-J(S)=F_S^\top A^{-1}B
 (C-B^\top A^{-1}B)^{-1}B^\top A^{-1}F_S.
 \label{eq:gating-gap}
\end{equation}
Equality holds exactly when $B^\top A^{-1}F_S=0$.
\end{proposition}
\begin{proof}
Positive definiteness of $R$ makes the diagonal blocks and both Schur
complements positive definite. Block inversion gives
$(R^{-1})_{SS}=A^{-1}+A^{-1}B(C-B^\top A^{-1}B)^{-1}B^\top A^{-1}$.
Congruence by $F_S$ proves the formula and positive semidefiniteness.
The positive definite middle matrix makes the resulting quadratic form zero
exactly under the stated condition. Empty and full selections give equality
directly.
\end{proof}
Equality covers independent errors, unions of complete independent
covariance blocks, and special sensitivity directions even when $B\ne0$.
In particular, correlation within an all-or-nothing sensor block causes no
mismatch when that block is independent of every other selection unit, so
the mismatch does not affect every correlated instance.

Gating has a statistical interpretation under a different experiment. If the
unselected raw noise values $\epsilon_{S^c}$ are supplied as
parameter-independent side information, the conditional mean of $Y_S$ is
$\mu_S+BC^{-1}\epsilon_{S^c}$ and the covariance is
$A-BC^{-1}B^\top$. Sensitivity remains $F_S$, giving
\eqref{eq:gated-information}. Equivalently, all channels can be observed
under the altered model $Y=D_S\mu(\theta)+\epsilon$.

Conditioning on actual unselected responses is different. At fixed
$Y_{S^c}=y_{S^c}$ the conditional mean is
$\mu_S+BC^{-1}(y_{S^c}-\mu_{S^c})$, so its sensitivity is
$F_S-BC^{-1}F_{S^c}$. Conditional information is consequently
\begin{equation}
 (F_S-BC^{-1}F_{S^c})^\top(A-BC^{-1}B^\top)^{-1}
 (F_S-BC^{-1}F_{S^c}).
 \label{eq:conditional-information}
\end{equation}
Adding $F_{S^c}^\top C^{-1}F_{S^c}$ gives full-response information.
Describing gating as ``conditioning on the other observations'' without
this sensitivity correction is therefore generally incorrect.

\begin{example}[Exact information and decision discrepancies]
\label{ex:gating-witnesses}
Let $J_0=0$, $p=1$, $R=\left(\begin{smallmatrix}1&3/4\\3/4&1\end{smallmatrix}\right)$,
and $F=(1,1)^\top$. Observing only the first response gives marginal
information $1$ and gated information $16/7$. Observing both gives $8/7$;
the first response conditional on the actual second has information $1/7$.
The gated score therefore decreases when the second response is selected.

The discrepancy can change an optimum. Select exactly one of three equal-cost
candidates with
\[
 R=\begin{pmatrix}1&0&3/4\\0&1&0\\3/4&0&1\end{pmatrix},\qquad
 F=(1,6/5,0)^\top,\quad J_0=0.
\]
The marginal singleton values are $(1,36/25,0)$ and the gated values are
$(16/7,36/25,0)$. Their optimizers select the second and first candidates,
respectively. The same positive-information ranking applies to scalar
log-determinant and inverse information; the third design is singular.
\end{example}
In contrast, marginal information is monotone under adding observations.
For disjoint coordinate sets $S,T$, block inversion gives
\begin{equation}
 \begin{split}
 I(S\cup T)=I(S)&+(F_T-R_{TS}R_{SS}^{-1}F_S)^\top\\
 &\cdot(R_{TT}-R_{TS}R_{SS}^{-1}R_{ST})^{-1}
 (F_T-R_{TS}R_{SS}^{-1}F_S).
 \end{split}
 \label{eq:information-increment}
\end{equation}
The increment is positive semidefinite; the empty-set case is interpreted
directly. This identity does not assert submodularity of a scalar criterion.

\subsection{Uniform quantitative bounds for a gated model}
\label{subsec:gating-bounds}
The next result specializes the classical inverse-compression form of the
operator Kantorovich inequality \citep[Eq.~(1)]{Moradi2019}; it is not a new
matrix inequality.

\begin{proposition}[Condition-dependent comparison]
\label{prop:kantorovich}
Suppose $0<m\le M$ and $mI\preceq R\preceq MI$. Set
\begin{equation}
 \kappa=M/m,\qquad \alpha=\frac{(M+m)^2}{4Mm}
 =\frac{(\kappa+1)^2}{4\kappa}.
 \label{eq:alpha}
\end{equation}
Every schedule satisfies
\begin{equation}
 J(S)\preceq J_{\gate}(S)\preceq\alpha J(S).
 \label{eq:uniform-gating}
\end{equation}
The two information matrices have the same kernel. If $J(S)\succ0$ and
$r_S=\rank(R_{S,S^c})$, then
\begin{equation}
 0\le\log\det J_{\gate}(S)-\log\det J(S)
 \le\min(p,r_S)\log\alpha.
 \label{eq:rank-gating}
\end{equation}
For an empty or full selection set $r_S=0$.
\end{proposition}
\begin{proof}
For $\lambda\in[m,M]$, $(\lambda-m)(\lambda-M)\le0$, hence
$\lambda+Mm/\lambda\le M+m$. Applying this inequality to an orthogonal
eigendecomposition gives $R+MmR^{-1}\preceq(M+m)I$. Compress to selected
coordinates and put $A=R_{SS}$, $K=(R^{-1})_{SS}$. Then
\[
 K\preceq\frac{(M+m)I-A}{Mm}\preceq\alpha A^{-1}.
\]
The second inequality follows by diagonalizing $A$ and using
$t(M+m-t)\le(M+m)^2/4$ for its positive eigenvalues. Congruence by $F_S$,
and $J_0\preceq\alpha J_0$, prove the upper comparison; the lower comparison
is Proposition~\ref{prop:gating}. The sandwich equates zero quadratic forms
and hence kernels.
For $J(S)\succ0$, the eigenvalues of
$J(S)^{-1/2}J_{\gate}(S)J(S)^{-1/2}$ lie in $[1,\alpha]$.
Equation~\eqref{eq:gating-gap} bounds the rank of its difference from the
identity by $r_S$. At most $\min(p,r_S)$ eigenvalues differ from one;
the logarithm of their product proves \eqref{eq:rank-gating}.
\end{proof}

For positive diagonal $D$, replacing $(R,F)$ by $(DRD,DF)$ leaves both
information matrices unchanged for every coordinate selection; so do
invertible block-diagonal transformations when entire blocks are selected
together. Scaling to unit marginal variances therefore avoids measuring
conditioning in arbitrary physical units. Spectral enclosures for the scaled
covariance used in a certificate must themselves be valid.

\begin{corollary}[Consequences for design optimization]
\label{cor:gated-design}
Under the assumptions of Proposition~\ref{prop:kantorovich}, with $\alpha$
as in \eqref{eq:alpha}, suppose $\calF$ contains a positive definite information
design. Let $S_G$
maximize $\log\det J_{\gate}$ and let $S_*$ maximize $\log\det J$ over
$\calF$. Then
\begin{equation}
 \log\det J(S_G)\ge\log\det J(S_*)-p\log\alpha,\qquad
 \left(\frac{\det J(S_G)}{\det J(S_*)}\right)^{1/p}\ge\alpha^{-1}.
 \label{eq:gated-design-guarantee}
\end{equation}
If $r_S\le r$ uniformly over feasible schedules, the additive loss improves
to $\min(p,r)\log\alpha$. A global gated optimizer of weighted trace
information achieves at least $1/\alpha$ of the true optimum of that same
criterion. A global gated A-optimal design has true trace-inverse objective
at most $\alpha$ times the true A-optimum, on the common positive definite
domain.
\end{corollary}
\begin{proof}
The D-optimality chain is
\[
 \log\det J(S_G)\ge\log\det J_{\gate}(S_G)-p\log\alpha
 \ge\log\det J_{\gate}(S_*)-p\log\alpha
 \ge\log\det J(S_*)-p\log\alpha.
\]
Use \eqref{eq:rank-gating} for the refinement. Taking the trace after
multiplication by $W\succeq0$ preserves the sandwich and proves the trace
statement by the same optimality chain. Inversion reverses order:
$\alpha^{-1}J(S)^{-1}\preceq J_{\gate}(S)^{-1}\preceq J(S)^{-1}$.
Take traces and compare the two minimizers to obtain the A-optimality claim.
\end{proof}
These guarantees require global optimization of the gated criterion. An
additive gated log-objective suboptimality $\eta\ge0$ adds $\eta$ to the
loss in \eqref{eq:gated-design-guarantee}. A heuristic solution or a solver
termination status alone does not establish global optimality.

\begin{example}[Sharpness]
\label{ex:gated-sharpness}
For $0<\rho<1$, select exactly one observation from
\[
 R=\begin{pmatrix}1&\rho&0\\\rho&1&0\\0&0&1\end{pmatrix},\qquad
 F=(1,1/2,b)^\top,\qquad J_0=0.
\]
The eigenvalue bounds $m=1-\rho$, $M=1+\rho$ give
$\alpha=(1-\rho^2)^{-1}$. True singleton information values are
$(1,1/4,b^2)$, and gated values are $(\alpha,\alpha/4,b^2)$.
For $1<b^2<\alpha$, gating selects the first candidate and the true objective
selects the third. The true efficiency $1/b^2$ approaches $1/\alpha$ as
$b^2\uparrow\alpha$. Every singleton information matrix is positive.
Rational $\rho,b$ approximate the limit, so the efficiency lower bound is
sharp as an infimum for rational data, one parameter, and three candidates;
equivalently, the additive log-information loss has supremum $\log\alpha$.
Letting $\rho\uparrow1$ rules out a positive efficiency guarantee independent of
covariance conditioning.
\end{example}
For a weak-correlation family $R=\Lambda+\varepsilon\Upsilon$, with fixed
positive diagonal $\Lambda$ and fixed symmetric zero-diagonal $\Upsilon$,
the cross block $B$ is $O(\varepsilon)$. Equation~\eqref{eq:gating-gap}
gives $O(\varepsilon^2)$ information error because the inverse blocks stay
bounded near zero. This recovers the first-order agreement of
\citet[Proposition~2]{Liu2016}. Weak correlation is sufficient
asymptotically but not necessary for equality at a special selection.

\subsection{A source comparison and independent time-block repair}
\label{subsec:source-comparison}
\citet{Wang2024} formulate measurement costs for all-or-nothing sensors and
individually charged samples. We compare with arXiv:2406.09557v1 and the
accepted manuscript OSTI 2447595 (the publisher's typeset article was not
retrieved). In the accepted manuscript, Eqs.~(8)--(10) form pair
contributions from full inverse-covariance entries, and Eq.~(11) gates them
by joint selection. The public implementation at commit
\texttt{430090e610446aab88328ce495ffb15b684c56c4} computes the full
covariance pseudoinverse before extracting entries or blocks
\citep{MeasurementOpt2025}; for our positive definite examples this is the
ordinary inverse in \eqref{eq:gated-information}. The mismatch occurs with
one measurement modality, so it does not depend on cross-modality counting
conventions in the displayed formulation.

A local Fisher approximation does not remove a mismatch that is already
present for an exactly linear Gaussian mean. Conversely, the argument does
not invalidate the source's independent-noise case, whose rotary-bed
covariance is diagonal, and does not assert that every correlated optimizer
changes. Agreement between optimization packages that reuse the same
coefficients checks their stated objective, not the likelihood.

The source's kinetics code uses eight independent time blocks. At each time,
order three static-cost measurements (SCMs) before three dynamic-cost
measurements (DCMs), with species A, B, C in each group. The code covariance is
\begin{equation}
 R_6=\begin{pmatrix}C_0&C_0/2\\C_0/2&C_0\end{pmatrix},\qquad
 C_0=\begin{pmatrix}1&1/10&1/10\\1/10&4&1/2\\1/10&1/2&8\end{pmatrix}.
 \label{eq:source-covariance}
\end{equation}
The leading principal minors of $C_0$ are $1$, $399/100$, and $791/25$,
so $C_0\succ0$. The other Kronecker factor has eigenvalues $1/2$ and $3/2$,
so $R_6\succ0$. Same-species cross-modality correlation is $1/2$.
The SCM--SCM block of $R_6^{-1}$ is $(4/3)C_0^{-1}$, so selecting all SCMs
and no DCMs gives gated information exactly $4/3$ times the actual data
information. We do not assert that this pattern is feasible at every
budget. In Table~S-1 of the inspected accepted manuscript (printed p.~S-2,
PDF p.~46), the A-DCM row, C-DCM column is $0.01$, whereas the transposed
entry is $0.1$. Equation~\eqref{eq:source-covariance} uses the symmetric
public-code covariance and does not treat that asymmetric table as a valid covariance model.

Independence across time allows a direct repair. For each time $t$ and subset
$P\subseteq\{1,\ldots,6\}$, precompute
\begin{equation}
 Q_{t,P}=F_{t,P}^\top(R_6)_{PP}^{-1}F_{t,P},\qquad
 Q_{t,\varnothing}=0.
 \label{eq:independent-pattern}
\end{equation}
Choose one pattern per time with binary $z_{t,P}$ and
$\sum_P z_{t,P}=1$. Then
\begin{equation}
 J=J_0+\sum_{t,P}z_{t,P}Q_{t,P}
 \label{eq:pattern-sum}
\end{equation}
is exactly the marginal information: reorder selected coordinates by time,
invert the resulting block-diagonal covariance, and sum its sensitivity
quadratic forms. SCM presence is tied across time, and installation,
sampling and budget restrictions act on pattern membership. There are at
most $64$ patterns per time before pruning. This is ordinary finite pattern
enumeration, not a new hull theorem. It does not justify independent temporal
chains when same-time cross-channel covariance is present, or complete-state
Markov formulas under arbitrary partial observation of a latent state.

\subsection{What a global certificate establishes}
\label{subsec:certificate-contract}
For $f(S)=\log\det J(S)$ and a feasible positive definite design
$\widehat S$, a certificate supplies valid numbers $\ell,U$ such that
\begin{equation}
 \ell\le f(\widehat S)\le f_*:=\max_{S\in\calF}f(S)\le U.
 \label{eq:certificate-contract}
\end{equation}
The achieved D-efficiency is at least $\exp(-(U-\ell)/p)$.
Numerical primal values and support prices become rigorous certificates only
after feasibility, arithmetic enclosures, and the relevant inequalities and
model inputs have been justified. ``Exact design'' in design theory means a
discrete design rather than a continuous weight measure; it does not itself
mean globally certified optimality or exact numerical arithmetic.

By Proposition~\ref{prop:gating}, a global upper bound $U_G$ on the gated
log optimum also bounds $f_*$. Re-evaluating a feasible schedule with its
selected covariance supplies the lower side of
\eqref{eq:certificate-contract}. The gated solver's own gap is not the
marginal gap, but a valid bound from an existing model remains useful after
its statistical objective is corrected.

Our rational certificates treat stored decimal model entries as exact
rational numbers and enclose logarithms with outward bounds. They certify
the optimization problem specified by those inputs, not that numerically
generated sensitivities or estimated covariances equal physical quantities.
Sensitivity verification, covariance assumptions, finite-scenario coverage
and arithmetic certification answer different questions and are reported
separately.
